% ============================================================
% sec/03baseline_literatura.tex
%
% OBLIGATORIO: un número objetivo con su fuente, declaración
% explícita de qué se busca igualar o superar, y las diferencias
% de protocolo que impiden una comparación justa.
%
% Si el conjunto de datos carece de literatura previa, el
% requisito se sustituye por revisar trabajos con conjuntos
% análogos del mismo dominio.
% ============================================================
\section{Baseline de Literatura}
\label{sec:baseline}

\subsection{Ausencia de literatura sobre el conjunto de datos}
\label{subsec:ausencialiteratura}

El conjunto de datos empleado en este trabajo no cuenta, hasta donde
alcanza la revisión realizada, con publicaciones arbitradas que reporten
resultados sobre él. En consecuencia, no existe un mejor resultado
publicado sobre este conjunto con la métrica seleccionada, y el requisito
se satisface mediante la revisión de trabajos sobre conjuntos análogos del
mismo dominio, justificando en qué medida resultan comparables.

\razon{Esta declaración debe abrir la sección y no ocultarse. El enunciado
prevé explícitamente la sustitución cuando el conjunto carece de literatura
previa, de modo que declararla es cumplir el requisito y no eludirlo.
Además, la ausencia se recoge como brecha en la
Sección~\ref{subsec:brechas}.}

\subsection{Origen del punto de referencia}
\label{subsec:origenbaseline}

\pc{Descripción del certamen de minería de datos sobre un juego de cartas
coleccionables del cual procede el número objetivo: en qué consistía la
tarea, qué características del estado se empleaban, y bajo qué protocolo se
evaluó. Corresponde al Eje A y lo aporta Daniel.}

\pc{Valor del resultado oficial de referencia y del mejor resultado
alcanzado, con la métrica correspondiente y la cita al trabajo que los
publica.}

\razon{La cifra debe entrar acompañada de su fuente arbitrada en la
bibliografía. Un número de literatura sin referencia citable incumple el
criterio de esta sección, que exige declarar el resultado con su fuente.}

\subsection{Número objetivo declarado}
\label{subsec:objetivobaseline}

\begin{quote}
\textbf{Referencia declarada:} \pc{valor y métrica}, reportado sobre
\pc{conjunto de datos y juego} \cite{pc}.
\end{quote}

Este trabajo declara dicho valor como \textbf{referencia de orden de
magnitud} y no como cifra a superar. La distinción es sustantiva y se
justifica en el apartado siguiente.

\razon{Declarar el número como cifra a batir cuando procede de otro juego
constituiría una comparación no legítima. Declararlo como orden de magnitud
cumple el requisito de la sección y a la vez respeta el criterio de
comparabilidad.}

\subsection{Comparabilidad del protocolo}
\label{subsec:comparabilidad}

La comparación entre el resultado de referencia y los resultados que este
trabajo obtenga se considera \textbf{únicamente indicativa}. Las siguientes
diferencias impiden una comparación directa:

\begin{itemize}
  \item \textbf{Juego distinto.} El punto de referencia procede de un juego
        de cartas coleccionables con reglas, recursos y condiciones de
        victoria distintas de las del Pok\'emon Trading Card Game.
  \item \textbf{Población de jugadores distinta.} \pc{Precisar si el
        conjunto de referencia procede de partidas entre personas o entre
        agentes automáticos.} El conjunto empleado en este trabajo procede
        de partidas entre agentes automáticos seleccionados por
        clasificación, condición que se detalla en la
        Sección~\ref{sec:analisistecnico}.
  \item \textbf{Protocolo de partición distinto.} \pc{Precisar el protocolo
        del certamen de referencia.} Este trabajo emplea partición agrupada
        por episodio, condición que reduce el desempeño medido respecto de
        una partición aleatoria, según se argumenta en la
        Sección~\ref{sec:metodologia}.
  \item \textbf{Representación del estado distinta.} \pc{Precisar qué
        características empleaba el certamen de referencia.}
\end{itemize}

\razon{Esta enumeración es la exigencia central del criterio de baseline.
Hodge \textit{et al.} ofrecen el precedente: publican una tabla comparativa
de dieciséis trabajos y advierten en nota al pie que la variación en
tamaño, composición y versión del juego obliga a tratar con cautela la
comparación de las exactitudes reportadas
\cite{hodge_winprediction_2021}.}

\noindent\textbf{Interpretación del número.} La referencia indica el orden
de magnitud del desempeño que la literatura considera alcanzable en tareas
de predicción de desenlace a partir del estado en juegos de cartas
coleccionables. No constituye un umbral de éxito ni de fracaso para este
trabajo.

\subsection{Referencias adicionales de orden de magnitud}
\label{subsec:referenciasadicionales}

Dado que la comparación directa no es posible, se recogen dos referencias
adicionales del dominio de juegos competitivos que permiten situar el orden
de magnitud esperable en tareas de predicción de desenlace a partir del
estado de juego.

Hodge \textit{et al.} reportan exactitudes en el rango de setenta a ochenta
por ciento sobre datos mixtos y profesionales de un videojuego de arena
multijugador, evaluando en el punto medio de la partida
\cite{hodge_winprediction_2021}. \pc{Verificar los valores exactos de la
Tabla III del artículo antes de incorporarlos.}

Xenopoulos \textit{et al.} reportan valores de pérdida logarítmica en torno
a cuatro décimas y errores de calibración esperados de un orden de
magnitud considerablemente menor dentro del entorno de entrenamiento
\cite{xenopoulos_winprob_2022}. \pc{Verificar los valores exactos de la
Tabla 1 del artículo antes de incorporarlos.}

\razon{Estas dos referencias se incluyen con su métrica original y sin
convertirla. Traducir una exactitud a un área bajo la curva, o una pérdida
logarítmica a una exactitud, produciría cifras que ninguno de los dos
artículos reporta.}

\decidir{Queda pendiente definir si el trabajo declara además un umbral
propio de éxito, independiente de la literatura, formulado sobre los
modelos de referencia internos descritos en la
Sección~\ref{subsec:planbaselines}. Un umbral de esa naturaleza sería
comparable de forma legítima, a diferencia del número de literatura. La
decisión afecta la redacción de la discusión final y conviene tomarla antes
de la Etapa 2.}